\section{The Evidence-Typed Data Model}
\label{sec:model}

\begin{figure}[tbp]
\centering
\resizebox{\textwidth}{!}{% Figure: core graph ontology, grouped into object families.
% Deliberately an OVERVIEW. 55 node labels and 81 relationship types are in use;
% drawing all of them produces a picture nobody can read, so each family is a group
% and carries one labelled edge from the textual spine. The full catalogue is
% Appendix A. Counts are from the fact freeze at 6bf220c.
\begin{tikzpicture}[
  font=\footnotesize,
  every node/.style={align=center},
  ent/.style={draw=black!45, rounded corners=2pt, inner sep=3.2pt, fill=white,
              minimum height=6.0mm, line width=0.45pt, font=\scriptsize},
  spine/.style={ent, fill=SkyBlue!18, draw=SkyBlue!65!black, font=\small},
  sub/.style={ent, fill=SkyBlue!8, draw=SkyBlue!45!black},
  attr/.style={ent, fill=BurntOrange!14, draw=BurntOrange!70!black},
  sem/.style={ent, fill=RoyalPurple!11, draw=RoyalPurple!60!black},
  reuse/.style={ent, fill=SeaGreen!14, draw=SeaGreen!60!black},
  rit/.style={ent, fill=Goldenrod!22, draw=Goldenrod!75!black},
  meta/.style={ent, fill=black!7, draw=black!45},
  grp/.style={draw=black!25, rounded corners=3pt, inner sep=6pt, dashed,
              line width=0.5pt},
  ttl/.style={font=\scriptsize\bfseries, text=black!60, inner sep=1pt},
  rel/.style={-{Latex[length=1.7mm,width=1.3mm]}, draw=black!55, line width=0.5pt},
  lbl/.style={font=\scriptsize\itshape, inner sep=1.5pt, fill=white, text=black!65},
]

% ============ textual spine, centre-left column ============
\node[spine] (work) at (0,0)      {\textbf{Work}\,\tiny(4)};
\node[spine] (pass) at (0,-1.6)   {\textbf{Passage}\,\tiny 22{,}537\\[-2pt]
                                   \tiny\code{:Mantra} 20{,}210};
\node[sub]   (tv)   at (-2.15,-3.45) {TextVersion\\[-2pt]\tiny 59{,}922};
\node[sub]   (tr)   at ( 0.45,-3.45) {Translation\\[-2pt]\tiny 18{,}427};

\draw[rel] (work) -- node[lbl] {\code{CONTAINS}} (pass);
\draw[rel] (pass.south) -- ++(0,-0.42) -| (tv.north);
\draw[rel] (pass.south) -- ++(0,-0.42) -| (tr.north);
\node[lbl, anchor=east] at (-0.16,-2.28) {\code{HAS\_TEXT\_VERSION}};
\node[lbl, anchor=west] at ( 0.16,-2.72) {\code{HAS\_TRANSLATION}};

% self-loop for the four cross-textual predicates
\draw[rel, dashed] (pass.west) to[out=160,in=200,looseness=5] (pass.south west);
\node[lbl, anchor=east, text width=28mm, align=right, fill=none] at (-1.45,-0.55)
  {\code{EXACT\_PARALLEL\_OF},\\[-2pt]\code{NEAR\_PARALLEL\_OF},\\[-2pt]
   \code{VARIANT\_OF},\\[-2pt]\code{REUSES\_TEXT\_FROM}};

\node[meta] (src) at (0,1.5) {Source \tiny(9)\ \ $+$ rights registry};
\draw[rel, draw=black!40] (src) -- node[lbl] {\scriptsize per-file rights} (work);

\begin{scope}[on background layer]
  \node[grp, fill=SkyBlue!5, fit=(work)(pass)(tv)(tr)] (gspine) {};
\end{scope}
\node[ttl, anchor=north west] at ($(gspine.north west)+(2pt,-1pt)$) {textual spine};

% ============ attribution family, right, top ============
\node[attr] (dev) at (5.05,0.35)  {\code{Devata} \tiny 214};
\node[attr] (rsi) at (7.30,0.35)  {\code{Rishi} \tiny 729};
\node[attr] (chn) at (4.95,-0.50) {\code{Chandas} \tiny 575};
\node[attr] (asc) at (8.15,-0.50) {\code{DevataAscription} \tiny 324};
\begin{scope}[on background layer]
  \node[grp, fill=BurntOrange!4, fit=(dev)(rsi)(chn)(asc)] (gattr) {};
\end{scope}
\node[ttl, anchor=south west] at ($(gattr.north west)+(2pt,1pt)$)
  {traditional attribution\ \ \scriptsize\mdseries(\code{HAS\_DEVATA},
   \code{HAS\_RISHI}, \code{HAS\_CHANDAS})};

% ============ textual reuse family ============
\node[reuse] (form) at (5.45,-2.2) {\code{Formula} \tiny 4{,}825};
\node[reuse] (fam)  at (7.95,-2.2) {\code{FormulaFamily} \tiny 720};
\node[reuse] (ppg)  at (6.6,-3.0)  {\code{PadaParallelGroup} \tiny 1{,}434};
\begin{scope}[on background layer]
  \node[grp, fill=SeaGreen!4, fit=(form)(fam)(ppg)] (greuse) {};
\end{scope}
\node[ttl, anchor=south west] at ($(greuse.north west)+(2pt,1pt)$)
  {textual reuse\ \ \scriptsize\mdseries(\code{USES\_FORMULA},
   \code{HAS\_PARALLEL\_PADA})};

% ============ semantic family ============
\node[sem] (sa)   at (5.70,-5.05) {\code{SemanticAssertion} \tiny 35{,}131};
\node[sem] (conc) at (4.95,-5.90)  {\code{Concept} \tiny 384};
\node[sem] (de)   at (7.35,-5.90) {\code{DomainEntity} \tiny 384};
\node[sem] (rf)   at (8.85,-5.05) {\code{RoleFiller} \tiny 2{,}052};
\node[sem] (lem)  at (6.6,-6.70)  {\code{Lemma} \tiny 10{,}031
                                   \ \scriptsize\itshape(internal)};
\begin{scope}[on background layer]
  \node[grp, fill=RoyalPurple!4, fit=(sa)(conc)(de)(rf)(lem)] (gsem) {};
\end{scope}
\node[ttl, anchor=south west] at ($(gsem.north west)+(2pt,1pt)$)
  {semantic\ \ \scriptsize\mdseries(\code{HAS\_SEMANTIC\_ASSERTION},
   \code{ABOUT\_CONCEPT}, \code{MENTIONS\_LEMMA})};

% ============ ritual family ============
\node[rit] (rite) at (5.15,-8.35) {\code{Ritual} \tiny 103};
\node[rit] (step) at (7.75,-8.35) {\code{RitualStep} \tiny 3{,}121};
\node[rit] (offr) at (6.6,-9.15)  {\code{Offering}, \code{Substance},
                                   \code{Object}, \code{Plant}};
\begin{scope}[on background layer]
  \node[grp, fill=Goldenrod!5, fit=(rite)(step)(offr)] (grit) {};
\end{scope}
\node[ttl, anchor=south west] at ($(grit.north west)+(2pt,1pt)$)
  {ritual \& material culture\ \ \scriptsize\mdseries(\code{USED\_FOR\_RITE})};

% ============ diagnostics, bottom left ============
\node[meta] (qa)  at (-0.95,-7.25) {\code{QAIssue} \tiny 915};
\node[meta] (qv)  at (-0.95,-8.05) {\code{QualityVerdict} \tiny 2{,}568};
\node[meta] (sch) at (-0.95,-8.85) {\code{ScholarlyDisagreement} \tiny 113};
\node[meta] (ic)  at (-0.95,-9.65) {\code{InterpretiveClaim} \tiny 6};
\begin{scope}[on background layer]
  \node[grp, fill=black!3, fit=(qa)(qv)(sch)(ic)] (gmeta) {};
\end{scope}
\node[ttl, anchor=south west, text width=40mm, align=left]
  at ($(gmeta.north west)+(2pt,1pt)$)
  {scholarship \& diagnostics\\[-1pt]
   \scriptsize\mdseries the first two are \code{:Internal}};

% ============ spine-to-family edges ============
\draw[rel] (pass.east) -- ($(gattr.west)+(0,0)$);
\draw[rel] (pass.east) -- ($(greuse.west)+(0,0.25)$);
\draw[rel] (pass.south east) to[out=-40,in=175] ($(gsem.west)+(0,0.35)$);
\draw[rel] (pass.south east) to[out=-62,in=172] ($(grit.west)+(0,0.2)$);
\draw[rel, draw=black!40] (pass.south west) to[out=-115,in=75] (gmeta.north);

% ============ evidence-layer key ============
\begin{scope}[shift={(-2.9,-10.9)}]
  \node[ttl, anchor=west] at (0,0) {Epistemic layer carried by \emph{every} edge:};
  \foreach \i/\c/\t in {0/{SkyBlue!45}/{\textsc{l1} 21.4\%},
                        1/{SkyBlue!75}/{\textsc{l2} 77.2\%},
                        2/{RoyalBlue!75}/{\textsc{l3} 0.17\%},
                        3/{RoyalBlue!60!black}/{\textsc{l4} 0.35\%}} {
    \node[draw=black!40, fill=\c, minimum width=3mm, minimum height=3mm,
          inner sep=0pt, line width=0.4pt] at ({5.15+\i*1.55},0) {};
    \node[font=\scriptsize, text=black!65, anchor=west]
      at ({5.3+\i*1.55},0) {\t};
  }
\end{scope}
\end{tikzpicture}
}
\caption{Core graph ontology, grouped by object family. An overview: 55 node labels
and 81 relationship types are in use, and the full catalogue is
\cref{sec:appendix-ontology}. The figure shows one labelled edge per family rather
than every predicate, because the alternative is unreadable. Counts are from the
fact freeze at \code{6bf220c}.}
\label{fig:ontology}
\end{figure}

\subsection{Four axes, deliberately not collapsed}

The object families are shown in \cref{fig:ontology}. Every relationship in
VedGraph carries, as stored properties, the answers to four different questions. They are orthogonal, and the model's central claim is that
collapsing any pair makes a well-formed scholarly question unaskable.

\paragraph{(1) How did this claim come to be?} The \emph{knowledge layer} is a
four-member closed vocabulary, described in the code as ``never collapsed, never
inferred from a score'':

% generated by figures-src/make_tables.py -- do not edit
\begin{table}[tb]
\centering
\caption{Relationships by epistemic layer, measured on the live graph at
\code{6bf220c}. The layer is a stored, indexed edge property, so this distribution is
a query rather than an estimate. Two pipeline generations wrote the field with and
without the \code{L}-prefix; the legacy spellings are folded here and reported
unfolded in the fact freeze.}
\label{tab:layers}
\small
\begin{tabularx}{\linewidth}{@{}l>{\raggedright\arraybackslash}Xrr@{}}
\toprule
Layer & Meaning & Edges & \% \\
\midrule
\Lone & \footnotesize a source states it & 109{,}298 & 21.39 \\
\Ltwo & \footnotesize a reproducible rule over stored text or over \vocab{l1} & 394{,}540 & 77.22 \\
\Lthree & \footnotesize a language model read evidence and proposed it & 890 & 0.17 \\
\Lfour & \footnotesize a reading held by this project or a named commentator & 1{,}787 & 0.35 \\
untyped (a recorded defect) & \footnotesize no layer recorded; see \cref{sec:threats} & 4{,}390 & 0.86 \\
\midrule
\textbf{Total} & & \textbf{510{,}905}
& \textbf{100.00} \\
\bottomrule
\end{tabularx}
\end{table}


\paragraph{(2) How should a reader grade it?} The \emph{quality tier} is a
reader-facing letter derived from the layer by a total bijection
($\textsc{l1}\!\to\!\textsc{a}$, $\textsc{l2}\!\to\!\textsc{b}$,
$\textsc{l3}\!\to\!\textsc{c}$, $\textsc{l4}\!\to\!\textsc{d}$). The module states the
rule as: tier is how a claim was reached, never how confident it sounds. The tier
exists separately from the layer because one further transition acts on it --- a
model-extracted claim that no reviewer has accepted stays in \textsc{l3} but drops to
\textsc{tier\,d}, since \textsc{tier\,c} is defined as model-extracted evidence that
survived review and a candidate has not had any.

\paragraph{(3) What is the claim's reach over the passage?} \emph{Attribution
precision} distinguishes a per-passage statement, a label \emph{inherited} from an
enclosing container, the passage's own wording, and edges that are not attributions at
all. This is the axis that matters most for counting. The Anukrama\d{n}\=\i\ names a
deity for a \iast{s\=ukta}; projecting that label onto each of its mantras is what
makes ``passages of Indra'' work at all, but a reader who cannot see the difference
reads tens of thousands of inherited labels as per-verse statements. The model records
the consequence: ``How many mantras invoke Indra?'' has at least four honest answers
--- strict source-stated, container-inherited, textual mention, and semantic
invocation --- and collapsing any pair of them makes one of the four unaskable. The
live distribution shows how unevenly the corpus behaves: every one of the
V\=ajasaneyi Sa\d{m}hit\=a's 2{,}240 seer attributions is source-stated per verse,
while every one of the Atharvaveda's 5{,}084 is a \iast{s\=ukta} label projected
downward, so a per-verse Atharvavedic seer claim cannot be made from this corpus at
all.

\paragraph{(4) Off which textual surface was the evidence read?} A rule fired over
Griffith's 1896 English is exactly as reproducible as a rule fired over the Sanskrit,
so the tier cannot separate them --- and the difference matters enormously, because
one measures a Victorian translator's word choice and the other measures the text.
The project discovered this the expensive way: 21{,}246 of 47{,}542
\code{ABOUT\_CONCEPT} edges had fired on an English keyword with no Sanskrit evidence
and carried an identical grade to edges matched on the verse itself. Those edges were
later withdrawn, and a separate \emph{evidence surface} axis was introduced so the
distinction is filterable rather than archaeological.

\subsection{A name collision, and what it cost}
\label{sec:collision}

The fourth axis is instructive as a failure. Two different enumerations in the
codebase are both called \code{EvidenceBasis}. One describes the textual surface the
evidence was read off; the other describes the kind of act that produced the claim.
Their value spaces are disjoint. Reading the stored graph property straight into the
second enumeration --- the natural thing to do, given the shared name --- mapped
\emph{every} attribution edge in the corpus to \code{UNKNOWN}: all 17{,}889 seer
edges, all 16{,}298 metre edges, all 10{,}558 deity edges. The failure was silent,
because \code{UNKNOWN} is a legitimate value.

The remedy is worth reporting because it generalises. Rather than renaming one
enumeration, the project added a test that enumerates the property's entire observed
value space in the live store and asserts it is a subset of the surface vocabulary.
The test's header records why the check is shaped that way: this project has twice
certified an absence by asking the wrong question, once against the wrong surface and
once with the wrong vocabulary. Grepping for the value you expect confirms your
expectation; enumerating the value space does not.

\subsection{The lifecycle: a model may write exactly one status}
\label{sec:model-policy}

The constraint that keeps \textsc{l3} small is not a review convention. It is encoded
as data. Assertion state is a three-member vocabulary --- candidate, accepted,
rejected --- and the set of trust classes permitted to be written as \emph{accepted}
without review is a frozen set containing exactly deterministic-derived and
source-explicit. The comment beside it is explicit that this is the policy expressed
as data rather than as a convention. The constructor raises if a model-derived record
is marked accepted, and raises again if a record claims model provenance without
naming a model.

The governing decision record, \mbox{ADR-014}, states the rule and three consequences
that are each enforced in code:

\begin{quote}\small
A model may write exactly one status: \code{CANDIDATE}. Every other status is written
by a deterministic stage or by a person. \dots{} Evidence is checkable or the claim is
discarded: a cited token the model was never shown is a fabricated citation, it is
decidable, and it is rejected. \dots{} The predicate vocabulary is closed, and the
refusals are named. \dots{} Confidence is a floor inside a predicate's rule and
nothing else: acceptance requires that the predicate has been \emph{unlocked} by a
measured gold precision of $\ge 95\,\%$; until then, a confidence of 1.0 is accepted at
the same rate as 0.4 --- not at all.
\end{quote}

Its stated rationale is the reason we treat confidence-thresholding as unsound in this
setting: the confidence score is produced by the same process that produced the claim,
so it carries no information independent of the claim, and thresholding on it selects
for fluency rather than for truth. The unlocked-predicate set is empty by
construction and remains empty; consequently no semantic assertion has ever been
auto-accepted.

Two refinements are worth noting. First, five predicates that a model might naturally
want --- \code{SYMBOLIZES}, \code{REPRESENTS}, \code{IS\_GOD\_OF}, \code{MEANS},
\code{CAUSES} --- are members of the enumeration \emph{so that a proposal naming one
is rejected as a policy decision with a recorded reason}, rather than failing schema
validation as an unknown string. The difference matters when reading a rejection
report. Second, the strongest review state the graph has ever recorded is
\code{MODEL\_ADJUDICATED}, glossed in the API as ``a model re-read the passage and
accepted this edge with a stated reason. This is the strongest review state in this
graph and it is not human review.'' The count of human-reviewed edges is zero, and the
product says so.

\begin{figure}[tbp]
\centering
\resizebox{0.93\textwidth}{!}{% Figure: the epistemic pipeline -- how a claim acquires, and keeps, its typing.
% Centre column: the state the claim is in. Left: who acts. Right: what travels with
% the claim. Bottom: the four terminal dispositions, reached by dedicated routing
% channels either side of the centre column so no connector crosses annotation text.
\begin{tikzpicture}[
  font=\footnotesize,
  every node/.style={align=center},
  st/.style={draw=black!50, rounded corners=2.5pt, inner sep=4pt, fill=white,
             minimum width=30mm, text width=30mm, minimum height=8mm,
             line width=0.5pt},
  gate/.style={st, draw=BurntOrange!70!black, fill=BurntOrange!10},
  term/.style={draw=black!45, rounded corners=2.5pt, inner sep=3.5pt, fill=black!4,
               minimum width=26mm, text width=26mm, minimum height=8mm,
               line width=0.45pt},
  acc/.style={term, fill=SeaGreen!12, draw=SeaGreen!60!black},
  flow/.style={-{Latex[length=2mm,width=1.6mm]}, draw=black!60, line width=0.6pt},
  route/.style={-{Latex[length=1.8mm,width=1.4mm]}, draw=black!40, line width=0.5pt},
  lft/.style={font=\scriptsize, text=black!60, align=right, text width=40mm,
              anchor=east},
  rgt/.style={font=\scriptsize, text=black!60, align=left, text width=56mm,
              anchor=west},
  hd/.style={font=\scriptsize\bfseries, text=black!55},
]

\def\L{-2.35}   % right edge of the left annotation column
\def\R{2.35}    % left edge of the right annotation column

\node[hd] at (-4.3,1.15) {actor};
\node[hd] at (0,1.15)    {state of the claim};
\node[hd] at (5.2,1.15)  {what travels with the claim};

\node[st]  (src)  at (0,0)     {\textbf{Source observation}};
\node[st]  (cand) at (0,-1.95) {\textbf{Candidate}};
\node[gate](val)  at (0,-3.95) {\textbf{Validation}\\[-1pt]
                                \scriptsize deterministic; 20 rejection codes};
\node[gate](adv)  at (0,-6.15) {\textbf{Adversarial gate}\\[-1pt]
                                \scriptsize independently implemented};
\node[st]  (grade)at (0,-8.25) {\textbf{Grading}\\[-1pt]
                                \scriptsize layer $\to$ tier, applied last};
\node[st]  (imp)  at (0,-10.3) {\textbf{Plan $\to$ dry run}\\[-1pt]
                                \textbf{$\to$ import $\to$ readback}};

\foreach \a/\b in {src/cand, cand/val, val/adv, adv/grade, grade/imp}
  \draw[flow] (\a) -- (\b);

% ---- left column: actors ----
\node[lft] at (\L,0)      {hashed snapshot; rights recorded per file};
\node[lft] at (\L,-1.95)  {\textbf{agent} or deterministic rule \emph{proposes}};
\node[lft] at (\L,-3.95)  {\textbf{deterministic code} only};
\node[lft] at (\L,-6.15)  {\textbf{a second agent} attacks, with a negative control};
\node[lft] at (\L,-8.25)  {\textbf{one contract}, applied graph-wide};
\node[lft] at (\L,-10.3)  {\textbf{owner} issues \textsc{go}; code verifies the delta};

% ---- right column: what travels ----
\node[rgt] at (\R,0)     {\code{source\_id}, \code{source\_artifact\_id},
                          \code{source\_locator}, \code{content\_sha256},
                          \code{rights\_status}};
\node[rgt] at (\R,-1.95) {\code{state=CANDIDATE} --- the only status a model may
                          write. Every cited token must be present in the packet the
                          model was shown.};
\node[rgt] at (\R,-3.95) {one code for one failure mode, e.g.\
                          \code{EVIDENCE\_TOKEN\_NOT\_IN\_PASSAGE},
                          \code{PREDICATE\_FORBIDDEN}};
\node[rgt] at (\R,-6.15) {\code{SURVIVED} or \code{DEFECT\_FOUND}. The failing run is
                          preserved and pinned by a test, so the finding cannot be
                          erased by regenerating only the passing run.};
\node[rgt] at (\R,-8.25) {\code{knowledge\_layer}, \code{quality\_tier},
                          \code{attribution\_precision}, \code{evidence\_basis},
                          \code{grade\_basis}};
\node[rgt] at (\R,-10.3) {rows \emph{sent} vs rows \emph{landed} --- the difference is
                          where the defects are};

% ---- terminal dispositions ----
\def\TY{-12.95}
\node[acc]  (t1) at (-4.35,\TY) {\textbf{Accepted}\\[-1pt]
                                 \scriptsize in the graph, typed};
\node[term] (t2) at (-1.45,\TY) {\textbf{Withheld}\\[-1pt]
                                 \scriptsize kept; reason typed};
\node[term] (t3) at ( 1.45,\TY) {\textbf{Rejected}\\[-1pt]
                                 \scriptsize deleted; reason kept};
\node[term] (t4) at ( 4.35,\TY) {\textbf{Unresolved}\\[-1pt]
                                 \scriptsize named, not guessed};
\node[hd, anchor=west] at (-6.6,-11.95) {terminal dispositions};

% routing channels: x=-1.85 / -2.1 on the left, x=1.85 / 2.1 on the right.
% The annotation columns start at +/-2.35, so nothing crosses text.
\draw[flow] (imp.south) -- ++(0,-0.7) -| (t1.north);
\draw[route] (adv.west)  -- (-1.85,-6.15) -- (-1.85,-12.25) -- (t2.north);
\draw[route] (val.east)  -- ( 1.85,-3.95) -- ( 1.85,-12.25) -- (t3.north);
\draw[route] (cand.east) -- ( 2.10,-1.95) -- ( 2.10,-12.55) -- (t4.north);

\node[font=\scriptsize\itshape, text=black!55, text width=13.6cm, align=left]
  at (0,-14.3)
  {A withheld claim is neither a rejected one nor a statement that the thing does not
   exist. Each withheld population has a \emph{closed} class space, and the number of
   rows carrying no disposition at all is itself gated at zero --- an untyped absence
   is one a reader cannot distinguish from a verified zero.};
\end{tikzpicture}
}
\caption{The epistemic pipeline: how a claim acquires its typing and what it can
terminally become. The centre column is the state of the claim, the left gutter names
the actor, and the right gutter names what travels with the claim at that point.
Only one of the four terminal dispositions puts a claim in the graph; the other three
are recorded rather than discarded, because an absence a reader cannot distinguish
from a verified zero is the most misleading thing this graph could publish.}
\label{fig:epistemic}
\end{figure}

\subsection{Typed absence}
\label{sec:typed-absence}

\Cref{fig:epistemic} sets out the whole lifecycle, from a source observation to one
of four terminal dispositions. The model treats an untyped absence as a defect. Four distinct mechanisms record
\emph{why} something is not present, each scoped to a layer: a reason string per
unmapped predicate, emitted into the load report alongside a separate count of
predicates that are neither mapped nor explicitly withheld (``not mapped'' and
``nobody looked'' are indistinguishable in a coverage number, and only one of them is
acceptable); a typed node property with a declared, closed class space, used for
S\=amavedic notation (\cref{sec:case-samaveda}); typed non-resolution rows in the
lexical layer, where 711 unresolved tokens each carry a linguistic reason rather than
a queue state; and edge-type-keyed withholding, where near-parallel pairs are held
back from a reuse claim with the reason recorded on the type.

The invariant that makes the S\=amavedic case work is published as a measurement
rather than asserted: the count of verses with no notation disposition at all must be
zero, because an untyped verse is one a reader cannot tell apart from a verified zero.

\subsection{The public boundary, and a label that does two jobs}

Product and internal knowledge share one database. Diagnostic nodes --- the
repository's assessments of its own passages --- carry a marker label so that a single
filter, written once, excludes them. The design note is that sprinkling an exclusion
across thirty queries would be the same policy enforced thirty times and wrong the
first time a new diagnostic label appeared.

The marker turns out to mark two different kinds of node, and the distinction is not
drawn by the label itself. It marks the diagnostic layer, which no researcher should
traverse into, and it marks \emph{sub-entities} --- text versions, translations,
lemmas, sources --- which must stay reachable, since a passage that could not reach
its own text would be useless. The broad form of the obvious leakage check therefore
reports 70{,}559 ``leaks'', every one of them legitimate, and an auditor running it
would file a false critical finding. The invariant is consequently written against the
diagnostic labels by name rather than against the marker.

A related caveat is structural and is worth stating for anyone building on such a
graph: the separation is query-side, not structural. An unfiltered one-hop traversal
from any passage returns tens of thousands of text-version and translation nodes. The
mitigation the project adopts is to serve a catalogue of vetted queries rather than
the database.

\subsection{Invariants}
\label{sec:invariants}

Twenty live invariants and twelve integrity gates run against the loaded store. The
design principle is stated in the invariant suite's header: each check is phrased as
the defect it detects rather than as the desired state, because each corresponds to a
defect this repository actually shipped. Representative members include: no edge
without a quality tier; no edge without a grade rationale; no unlabelled node; no
product node reaching a diagnostic node; no work without a scope statement; no
attribution-precision value outside the declared enumeration; no predicate carrying
more than one run identifier, since that would mean part of the graph is reproducible
from no committed artefact.

Two features of this suite bear on method. First, one gate had been rewritten twice
and both earlier versions were untestable --- they computed the set of declared
relationship types as the union of the ontology \emph{and the graph}, then asked which
graph types were missing from it. Nothing can be. The gate reported zero undeclared
types for an entire wave while eleven predicates went unclassified. The code now
carries the lesson as a comment: the data may not declare itself valid. Second, a
universal-quantifier invariant over attribution precision could only be written safely
once the vocabulary had a member meaning ``this edge is not an attribution'' ---
an invariant asserting that every edge carries a property must first be able to say
what the property means on every edge.
